% The spoken body of the homily for the feast of All Saints. Continuous
% preaching, ready to read aloud: no heading, no direction to the preacher,
% no citation inside the speech. The loci, the relation to the study's
% interpretations, the audience, the word count and the pace estimate stand
% in the terminal note, which no one reads out. The homily is preached after
% the Gospel and before the Credo, so the chants and prayers before the
% Gospel are spoken of as heard and those after the Credo as still to come.
% The blank lines are the paragraph breaks a speaker recovers his place by,
% and the vertical spaces divide the speech into its parts. It opens on the
% feast and on the Lesson's vision, heard number and seen crowd, and sets
% Bede's reading of the number beside Augustine's mass of the saints. It
% then asks the way into that crowd and answers from the Alleluia's
% invitation and the Gospel's eight blessings, with Augustine's ordered way
% and Chrysostom's golden chain and blessing without command. It hears the
% Gradual's promise against the saints' hunger and death, with Augustine's
% objector, Basil's perfect good and Aquinas's physician, and carries the
% answer to the Offertory's just who seemed to die, with Schuster's hands,
% and to the Secret over the gifts. It sets the present tense of the first
% and eighth blessings with Aquinas's way and homeland, and joins the
% crowned to the journeying in the Collect and the Postcommunion. It takes
% the commemorated Sunday's Collect as the prayer of the first step, poverty
% of spirit, and draws the one response from it; and it closes at the
% Communion, where the Church sings the three last beatitudes over the
% Sacrament, on the promise that the clean of heart shall see God.
% The honorific is tied to the name so that "St" never ends a line.
\raggedcolumns

This morning the Church keeps the feast of all the saints in one
celebration. The Mass opened with the Church's own words, not Scripture's:
``Let us all rejoice in the Lord, celebrating a feast in honor of all the
saints, in whose solemnity the angels rejoice and join in praising the Son
of God.'' William Durandus, who wrote the medieval Church's great book on
the liturgy, gives the plain reason for keeping them together: the saints
are too many for each to have his own day. So today the Church names no one
of them, and rejoices in them all.

The Lesson then gave us St~John's vision. An angel comes up from the east,
carrying the sign of the living God, and orders the angels who have power to
harm the earth to hold back until the servants of God are signed on their
foreheads. John hears the number of the signed: a hundred and forty-four
thousand, twelve thousand from each of the twelve tribes of Israel. Then he
looks, and he sees ``a great multitude, which no man could number, of all
nations and tribes and peoples and tongues,'' standing before the throne and
in sight of the Lamb, robed in white, with palms in their hands.

The Lesson gives an exact number, and then a crowd past counting. Taken as a
count, the number shuts the crowd out; taken as a figure, it seems to be set
aside for the crowd. St~Bede, the English monk who expounded this book verse
by verse, keeps both. ``By this definite number,'' he says, ``is signified
the innumerable multitude of the whole Church,'' gathered from the
patriarchs onward, by birth and by faith --- for, as St~Paul says, ``if ye
be Christ's, then are ye Abraham's seed.'' The seal on the forehead Bede
reads as the sign of the cross, the faith confessed openly where all can
see. The number says that God loses none of his own; the crowd says how many
his own are. St~Augustine, preaching on the few and the many, points to this
vision and says, ``This is the mass of the saints.'' The saved look few
beside the world. In God's sight they cannot be counted.

Bede also says what their robes and palms mean. The white robe is baptism,
the grace in which they were cleansed; the palms are ``the triumph of the
Cross,'' for ``in Christ they have overcome the world.'' And the angels of
the opening song are in the vision too: they fall on their faces before the
throne and give God a sevenfold praise.

\medskip

That crowd was gathered by an invitation, and we heard it sung before the
Gospel: ``Come to me all you that labor and are burdened, and I will refresh
you.'' St~John Chrysostom remarks that the Lord calls ``not this or that
person, but all that are in anxiety, in sorrows, in sins.'' And the Gospel
we have just heard shows the way by which those who answer come. Jesus sees
the multitudes, goes up the mountain, sits down as the teacher, and opens
his mouth. St~Augustine calls this sermon ``a perfect standard of the
Christian life,'' and he reads the eight blessings as one way in order. It
begins with the poor in spirit --- the humble, he says, who fear God ---
and it climbs through meekness, through mourning, through hunger and thirst
for justice, through mercy and purity of heart, to the peacemakers; and the
eighth blessing, of those persecuted for justice' sake, returns to the first
and completes it. Chrysostom calls the eight ``a sort of golden chain,''
each link drawing the next, and he points out how the Lord speaks: by
blessing, not by command, so that ``though thou be a slave, a beggar, in
poverty, a stranger, unlearned, there is nothing to hinder thee from being
blessed.'' The way is open to everyone who will take its first step, and
that is how a numbered Church becomes a crowd no man can count.

\medskip

The way is hard, and the Mass does not hide it. The Gradual sang: ``Fear the
Lord, all ye his saints: for there is no want to them that fear him.'' St~Augustine
hears the objection any of us could raise. We have seen wicked men
die rich; and the man who prays every day, and has nothing, says, ``the
Scripture has deceived me.'' The saints themselves went hungry, and some
were put to death. St~Basil answers that the perfect good is God himself:
those who seek the Lord lack no good, because they have him. St~Thomas
Aquinas adds that God is both nourisher and physician. A physician withholds
food from a sick man because it serves his healing; so God, as it serves our
salvation, sometimes sends want and sometimes gives plenty --- and he comes
to the help of those who seek him in their need.

The Mass presses the same answer as far as death. At the Offertory the
Church will sing: ``The souls of the just are in the hand of God, and the
torment of malice shall not touch them: in the sight of the unwise they
seemed to die, but they are in peace.'' Blessed Ildefonso Schuster, who
wrote on this Mass in the last century, says it in one sentence: ``The
persecutor believes that the lives of the martyrs are in his hands, but they
are in the hands of God.'' Then the priest, speaking quietly over the gifts
in the prayer called the Secret, will ask that our offering be pleasing to
God for the honor of all the just --- his word for the saints in that prayer
--- and that it be made profitable to us who offer it. In a few moments the
sacrifice of the Lamb, in whose sight the white-robed crowd is standing,
will be offered on this altar; the palms in their hands are the triumph of
his cross, and that same offering is made present here for us.

\medskip

Two of the eight blessings speak in the present tense, the first and the
last: ``theirs \emph{is} the kingdom of heaven.'' St~Thomas explains that
the rewards are had perfectly only in the homeland, and in their beginning
on the way. The saints have them whole; we have them begun. The prayers of
this Mass hold the two together. In the Collect we said that God has granted
us to venerate the merits of all his saints in one celebration, and we asked
that, with our intercessors so multiplied, he would grant the abundance of
his mercy. After Communion we will ask that we may always rejoice in
venerating all the saints and be ``defended by their unceasing prayers.''
The crowd before the throne is not closed against us. They pray for us from
there, and their prayer does not cease.

Because the feast falls this year on a Sunday, the priest said a second
Collect after the feast's own: the Collect of the Twenty-third Sunday after
Pentecost, the Sunday this feast outranks. It asks God to loose the sins of
his people and to deliver us from ``the trammels of our sins, in which,
through our frailty, we have become entangled.'' The same Mass that shows us
the crowned saints lets us ask pardon for our frailty; and that asking is
where the Gospel's way begins, because the first blessing is for the poor in
spirit, the man who knows his poverty and tells it. So take the Church's
prayer home and make it your own: tonight, or when you next prepare for
confession, tell God plainly where your frailty has bound you, and ask him
to loose it. That asking is the first step of the way the saints walked.

\medskip

When you come forward in a little while, the Church will sing three of these
blessings again: ``Blessed are the clean of heart: for they shall see God.
Blessed are the peacemakers: for they shall be called the children of God.
Blessed are they that suffer persecution for justice' sake: for theirs is
the kingdom of heaven.'' She sings them at the moment you receive him,
because in that Sacrament the end of the way is pledged to those still
walking it: the vision of God, the sonship, the kingdom. The saints in the
vision were once where you are. They heard the same invitation, they began
at the same first blessing, and now they stand before the throne and pray
for you without ceasing. The blessing sung at your Communion is the way's
end, and it is promised to you: the clean of heart shall see God.
